\documentclass[11pt,letterpaper]{article}
\usepackage[T1]{fontenc}
\usepackage{lmodern,microtype,amsmath,amssymb,amsthm,mathtools,booktabs,tabularx,array}
\usepackage[left=.89in,right=.89in,top=.89in,bottom=.87in]{geometry}
\usepackage[dvipsnames]{xcolor}
\usepackage{enumitem,fancyhdr}
\usepackage[breaklinks,colorlinks,urlcolor=MidnightBlue,linkcolor=MidnightBlue,citecolor=MidnightBlue]{hyperref}
\usepackage{url}
\urlstyle{tt}
\newtheorem{theorem}{Theorem}[section]
\newtheorem{lemma}[theorem]{Lemma}
\newtheorem{corollary}[theorem]{Corollary}
\newcommand{\Z}{\mathbb Z}
\newcommand{\Q}{\mathbb Q}
\newcommand{\C}{\mathbb C}
\newcommand{\cO}{\mathcal O}
\newcommand{\src}[3]{\href{https://github.com/VladimirReshetnikov/ProveIt/blob/89d12dd40424f1217fd818d84877b3da52396a3b/Analysis/GeneralizedRiemannHypothesis/Lean/grh/src/#1.lean\#L#2}{#3}}
\newcommand{\docref}[2]{\href{https://github.com/VladimirReshetnikov/ProveIt/tree/89d12dd40424f1217fd818d84877b3da52396a3b/Analysis/GeneralizedRiemannHypothesis/docs/#1}{#2}}
\newcommand{\db}[2]{\src{New/Congruences/FirstOrder/DenominatorBound}{#1}{#2}}
\newcommand{\lv}[2]{\src{New/OneLift/LiftVolumeHeight}{#1}{#2}}
\newcommand{\exis}[2]{\src{New/Reduction/Data/Existence}{#1}{#2}}
\DeclareMathOperator{\ord}{ord}
\setlist{topsep=3pt,itemsep=2pt}
\setlength\headheight{14pt}
\pagestyle{fancy}\fancyhf{}\lhead{\small Three simplifications of the theta-lift proof}\rhead{\small Research note}\cfoot{\thepage}
\setlength{\emergencystretch}{2em}
\hypersetup{pdftitle={Three local simplifications of the theta-lift proof of Dirichlet GRH},pdfauthor={OpenAI}}
\begin{document}
\begin{center}
{\LARGE\bfseries Three local simplifications of the\\[4pt] theta-lift proof of Dirichlet GRH}\par
\vspace{7pt}
{\large Denominator clearing, finite-field types, and sparse places}\par
\vspace{6pt}
{\small November 27, 2026}
\end{center}
\begin{abstract}
The determinant lower bound in the Alpöge--Claude proof passes through an operator preserving Hecke eigenspaces. Its application only requires an operator that clears prescribed projections and has a large local determinant. In an integral basis adapted to the flag of lift generators, taking one row from each available operator gives such a determinant exactly. The same integral basis argument shortens the estimate for the determinant of a compressed lift. A second substitution enlarges the family of finite-field labels with a surviving Frobenius orbit and a common norm of sufficient order. At least $q^r/(2r)$ degree-$r$ labels qualify, giving a witness-loss bound $e^{-f/(2r)}$. A third substitution constructs the sparse infinite places from the graph of a quadratic polynomial and removes the choice of a large sphere. Each substitution has its own insertion point in the readable Lean development.
\end{abstract}

\section{What the comparison proves}
At the center of the argument is a family of forms whose covolume is estimated twice. The theta derivatives make the family short after the known prime-power divisibilities have been divided out (Theorem A). A second-moment calculation provides sufficiently many independent forms among a controlled number of translates (Theorem B). The lower bound for the covolume of integral forms comes from a chain of theta lifts, special values, and the product formula (Theorem C). In that lower bound local denominators must be cleared without spending too much of the determinant. The first proposal changes this particular operation.

In the source notation, $D$ is admissible data, $m$ is the number of lifts, $f$ the number of infinite places, and $a=xb$ the large derivative order. If $d_0$ is the target dimension, $n$ is the number of independent test forms, $V$ bounds their logarithmic squared covolume, and $\delta$ bounds their index, the two sides of the comparison have the form
\[
 M_0-\text{error}\ \le\ \frac{V+2\delta}{d_0}
 \ \le\ M_0-\frac{n}{d_0}m\,\mathrm{rem}+\text{error}.
\]
Here $M_0$ is the same main amount on both sides. At the published parameters $n\ge d_0/2$, and the comparison gives
\[
                  m\,\mathrm{rem}\le f a(\log b)^5.
\]
The prime remainder contains a near cutoff at $x$ and a distant cutoff at $x_1$. The far term has a classical prime-number-theorem estimate. At the near cutoff $m\ge\sqrt x$ turns the arithmetic bound into a one-sided estimate at scale $x^{-1/2+\varepsilon}$. The Mellin transform detects off-line zeros. The choice of a quadratic twist prevents its zeros from cancelling a proposed pole, and Landau's argument for a positive transform excludes the pole. This gives the half-plane theorem for Dirichlet characters. Specialization to zeta and the functional equation give RH. The changes below reuse the inputs that establish these analytic steps. See the \src{New/Reduction/MainInequality}{87}{covolume inequality}, the \src{New/Reduction/OneSidedBound}{200}{one-sided remainder bound}, the \src{New/Reduction/OneSidedBound}{239}{half-plane theorem}, and its \src{Main}{5}{specialization to RH}.

\begin{center}\small
\begin{tabularx}{\linewidth}{@{}p{.23\linewidth}XX@{}}\toprule
Input to the proof & Present argument & Proposed replacement\\\midrule
Denominators in Theorem C & Choose a simultaneous scalar combination over an infinite residue field. & Select rows in a basis adapted to the integral lattice and the flag.\\[3pt]
Finite-field types used in B and C & Demand that a witness root generate the whole multiplicative group. & Test its Frobenius orbit and the prime factors needed in its norm.\\[3pt]
Admissible data for the derivative bound & Find a well-populated sphere among digit vectors. & Write the squared norm into the top two digits of an explicit exponent.\\\bottomrule
\end{tabularx}\end{center}
The first change has the smallest effect on the rest of the development. It preserves the existing denominator interface and the original projector. The label change requires transporting a larger projector through the proof. The digit construction can be used with either choice of projector.

\section{Clearing the projection one row at a time}\label{sec:flag}
Fix a prime $\ell$ and the isomorphism $\iota:\C\simeq\overline{\Q}_\ell$ used for this local estimate. Write $K$ for the coefficient field with this absolute value and $\cO=\{x:|x|\le1\}$. Write $P_Z$ for the orthogonal projection onto the lift space $Z$. It is defined before applying $\iota$ to coordinates. Let $T$ be the set of $P_Zv$ for integral forms $v$ orthogonal to the false lifts. This restriction is essential in the one-generator theorem. In the following argument $T$ can be any subset.

The published one-generator theorem supplies an operator for each member $z_i$ of a basis of lift generators. On $Z$ the operators $G_i$ are diagonal in that common basis, every $G_i$ sends $T$ into the integral forms in $Z$, and the eigenvalue of $G_i$ on $z_i$ has a lower bound in absolute value. In the Lean theorem $G_i$ initially takes values in smooth forms. The identification on the generator basis that restricts its values to $Z$ is already made in the \db{1191}{proof of the operator bound}. We keep this identification. The existing proof next finds a linear combination serving all basis vectors, proves that it preserves Hecke eigenspaces, and translates its eigenvalue bounds into a determinant bound. The caller discards the Hecke-preservation field. Its actual output also includes integral pairings after multiplication by a stabilizer factor. See the \db{1030}{one-generator result}, the \db{1279}{caller}, and the \src{New/OneLift/TestedForms}{413}{definition of the denominator bound}.

\subsection{An integral basis compatible with a flag}
Suppose $E:Z\hookrightarrow K^I$ is an injective linear map for a finite set $I$. Give $Z$ the lattice $L=E^{-1}(\cO^I)$. Let $z_1,\dots,z_d$ be any ordered $K$-basis and put $F_j=\langle z_1,\ldots,z_j\rangle_K$, with $F_0=0$.
\begin{lemma}\label{lem:pivot}
There are $b_j\in F_j$ and distinct $t_j\in I$ such that $b_1,\ldots,b_j$ span $F_j$, $b_j\in L$, $E(b_j)_{t_j}=1$, and $E(b_j)_{t_h}=0$ for $h<j$. They give
\[
             L\cap F_j=\bigoplus_{i=1}^j\cO b_i\qquad(0\le j\le d).
\]
\end{lemma}
\begin{proof}
For the $j$th step start with $u=z_j$. In order $h=1,\ldots,j-1$, replace $u$ by $u-E(u)_{t_h}b_h$. The previous pivot coordinates stay zero. The resulting $u$ has the class of $z_j$ modulo $F_{j-1}$, so $E(u)\ne0$. Choose a coordinate $t_j$ of maximum absolute value. It is a new pivot. Then $b_j=u/E(u)_{t_j}$ has all its coordinates in $\cO$ and has the claimed pivot values. The change of basis preserves each $F_j$.

If $v\in L\cap F_j$, apply the same ordered subtractions at $t_1,\ldots,t_j$. Each coefficient lies in $\cO$, since each intermediate remainder is still integral. The final remainder is zero: the evaluation matrix on the pivots of $b_1,\ldots,b_j$ is triangular with diagonal entries $1$. The opposite inclusion follows from $b_i\in L$. The proof uses only finite pivot searches and the ultrametric inequality, so it applies to nondiscrete valuation rings.
\end{proof}
\noindent
In the source, the finite coordinate map evaluates the form at representatives of the adelic double classes that are integral above $\ell$, then takes its Weyl coordinates. Vanishing on these representatives implies vanishing of the form. At such a representative the gauge bound is precisely the coordinate bound. Transformation by the rational and level subgroups transports that gauge over the full class, since the level character is a unit. These facts identify the source's integral forms on $Z$ with $L$ above. Choosing the integral representatives matters here. The ingredients appear together in \lv{124}{the normalization lemma}, with the \src{Lib/Adelic/AutomorphicForms/PAdicGauge}{150}{coordinate formula for the gauge} and \src{Lib/Adelic/AutomorphicForms/PAdicGauge}{232}{existence of representatives}.

\subsection{The row construction and its determinant}
\begin{theorem}\label{thm:rows}
For each $j=1,\ldots,d$, suppose $G_j:Z\to Z$ preserves the complete flag $F_\bullet$ and $G_j(T)\subseteq L$. Write $\gamma_j$ for its scalar on $F_j/F_{j-1}$. There exists an endomorphism $H$ of $Z$ satisfying
\[
                        H(T)\subseteq L,\qquad \det H=\prod_{j=1}^d\gamma_j .
\]
\end{theorem}
\begin{proof}
Take the basis from Lemma~\ref{lem:pivot} and its coordinate functionals $b_j^*$. Define
\[
                         Hx=\sum_{j=1}^d b_j^*(G_jx)b_j .
\]
For $x\in T$ each coefficient is integral by the lemma. For $x\in F_k$, all terms with $j>k$ vanish. In particular the $k$th coordinate of $Hb_k$ is $b_k^*(G_kb_k)=\gamma_k$. The matrix of $H$ is upper triangular with this diagonal, proving the identity. For $d=0$ the empty product is $1$.
\end{proof}

Equivalently, in the integral basis $H$ takes row $j$ from the matrix of $G_j$. Different operators can supply different rows even when some $G_j$ are singular. The theorem allows arbitrary residue fields and does not require simultaneous diagonalizability. In the application the lift eigenbasis supplies the flag. Its $j$th quotient sees exactly the eigenvalue on the original $z_j$. If candidates have been obtained for every direction, one may also choose for each quotient the candidate of largest absolute value there.

\paragraph{Sharpness with the stated data.}
Let $K$ have a uniformizer $\pi$, $Z=K^d$, $L=\cO^d$, and $T=\bigoplus_j\pi^{-s_j}\cO b_j$ with $s_j\ge0$ integers. Take $G_j$ to be $\pi^{s_j}$ times the $j$th coordinate projection. Any map that sends $T$ into $L$ has its $j$th column divisible by $\pi^{s_j}$. Thus every clearing determinant is divisible by $\pi^{\sum s_j}$. The construction attains that valuation. This example also explains the relevance of a product of individual bounds when they differ between generators.

\subsection{Restoring the stabilizer factor and integral pairings}
Use $v_\ell(\ell)=1$. Suppose the published inputs give $|\gamma_j|_\ell\ge \ell^{-B_j}$, and put $c=D.\mathrm{chosenGroupOrder}$. Taking $\Lambda=cH$ gives
\[
    |\det\Lambda|_\ell\ge
           \ell^{-d v_\ell(c)-\sum_j B_j}.
\]
For the uniform source bound $B_j=B$, the existing numerical inequality is $B+v_\ell(c)\le r'\beta_\ell$, with $r'$ the rank in the descent pair and $\beta_\ell=\mathrm{clearingBound}(D,\ell)$. Hence $|\det\Lambda|_\ell\ge\exp[-d r'\beta_\ell\log\ell]$. The $c^2$ already absorbed into the single-generator maps stays there. The factor $c$ here is the additional factor appearing in the final pairing estimate.

For an admissible $v$, the form $w=H P_Zv$ is integral at $\iota$. The source's stabilizer lemma gives $|\iota(c\langle u,w\rangle)|\le1$ when $u$ is integral at all places over $\ell$. Since $c$ is a rational integer, $c w$ is integral and $\langle u,cw\rangle=c\langle u,w\rangle$. These are both clauses required of $\Lambda$. When $Z=0$ the original identity solution on the zero space applies. When $\dim Z=1$ a single generator operator suffices. Thus the same proof serves both consumers, for the full lift and for the rank-one auxiliary lift. See the \src{Lib/Adelic/AutomorphicForms/PAdicGauge}{750}{pairing estimate}, the \db{1276}{local budget}, and the \src{New/Congruences/FitTwistsExist}{415}{full-rank} and \src{New/Congruences/FitTwistsExist}{444}{rank-one} applications.

\paragraph{A second use of the same basis lemma.}
At a prime $\ell$ prime to the twist modulus norm, take the comparison operator $U$ and set $S=P_ZU|_Z$. The source shows that $U$ sends every integral lift into an integral form orthogonal to false lifts. Consequently $(\Lambda S)L\subseteq L$. In an integral basis its matrix has entries in $\cO$, so $|\det(\Lambda S)|\le1$ by the determinant formula. The lower bound on $|\det\Lambda|$ immediately gives the required upper bound on $|\det S|$. This replaces the separate normalization of eigenvectors at the \lv{652}{compressed-determinant estimate}. It works for any clearing operator satisfying the same contract.

\subsection{Why the choice of basis matters}
Over $\Q_2$, with ordinary dot product on $\Q_2^3$, take $b_1=e_1$, $b_2=e_2+e_3$, and $Z=\langle b_1,b_2\rangle$. Its integral lattice has basis $b_1,b_2$. Orthogonal projection sends $e_2$ to $b_2/2$. Use a generator basis $z_1=b_1$, $z_2=b_1+2b_2$ and the two candidate operators $4I$ and $2I$. Each clears the projection of every integral vector. Prescribing eigenvalues $(4,2)$ in the $z$-basis sends $b_2/2$ to $-b_1/2+b_2$, which fails integrality. Taking the rows in the integral $b$-basis gives $H=\operatorname{diag}_b(4,2)$, $H(b_2/2)=b_2$ and $\det H=8$. This $H$ moves the line $Kz_2$ and so illustrates precisely the extra freedom used by the theorem.

The guide also exhibits $G=\operatorname{diag}(1,0,1)$ and $G'=\operatorname{diag}(0,1,1)$. An integral scalar combination over $\Z_2$ has diagonal $(u,v,u+v)$, which always contains an even entry. Taking rows from $G,G',G$ in the standard integral basis gives $I_3$ over $\Z_2$. The output lies outside the integral linear span of the candidates. For an application that genuinely needs a scalar combination or Hecke commutation, the published stronger theorem remains valuable.

\paragraph{Implementation boundary.}
A Lean implementation can first extract the finite evaluation equivalence for arbitrary vectors from the normalization lemma, then prove Lemma~\ref{lem:pivot} by induction on an ordered basis. It can reuse the construction of the generator basis and the restriction from $G_j$ to $Z$, and apply Theorem~\ref{thm:rows} directly in \db{1246}{the final clearing theorem}. This bypasses the \src{New/Congruences/FirstOrder/HeckeStrings}{59}{simultaneous specialization} on this branch and the \db{1219}{eigenvalue-to-determinant estimate}. Other declarations in these modules are still used elsewhere. Any reduction of source lines or build time depends on the eventual implementation.

\section{A denser choice of finite-field types}
\label{sec:denser-labels}

The projector in the guide's Section~10 has two jobs. At every chosen place its label must avoid certain linear roots. Above each of the two level primes at least one label must also furnish a cuspidal test and a tame-inertia witness for the denominator argument. The original construction uses a primitive root of the entire extension field for the second job. A positive proportion of the elements of the extension field suffice. This removes Euler's totient from the estimate and improves the sufficient number of places for this step from
\[
  18r(2r\log q+1)\log b \quad\text{to}\quad 18r\log b.
\]
Here $q$ is a level prime, $r$ is the target rank, and $b$ is the parameter in the bad-type bound. Other uses of the number of places in the global construction have their own constraints.

\paragraph{Finite-field theorem.}
Let $q\ge19$ be an odd prime power with $q\equiv1\pmod9$, let $r=3^a$ with $a\ge1$, and suppose $N^2\mid q-1$. Assume $v_p(N)\ge6$ for every prime $p\ne3$ dividing $N$. In the source $N=q_2(s)=(q_0d)^{s+6}$. Put
\[
 R=q^r-1,\qquad A=q^{r/3},\qquad C=A^2+A+1=3H,\qquad e=R/(q-1).
\]
Fix a generator $g$ of $\mathbb F_{q^r}^{\times}$. Select the minimal polynomials of the elements $z=g^k$ for which
\begin{equation}
  3\nmid k,\qquad H\nmid k,\qquad p^6\nmid k\quad(p\mid N,\ p\ne3).
  \label{eq:label-selection}
\end{equation}
There are exactly
\begin{equation}
 \frac{R}{r}\,\frac23\left(1-\frac1H\right)
              \prod_{\substack{p\mid N\\p\ne3}}(1-p^{-6})
 \label{eq:label-count}
\end{equation}
such polynomials, at least $q^r/(2r)$. Each is an irreducible label of degree $r$. For every $N$th root of unity $h\in\mathbb F_q$ and every selected root $z$, all the conjugates $\xi$ of $hz$ satisfy the following conclusions.
\begin{enumerate}
\item $\xi^{q^r}=\xi$. For every prime $\ell$ and integer $u\ge0$, the elements $(\xi^{q^i})^{\ell^u}$, $0\le i<r$, are pairwise distinct.
\item Their common $e$th power $y$ belongs to $\mathbb F_q^{\times}$ and $N\mid\operatorname{ord}(y)$. In particular $y^x=1$ implies $N\mid jx$ for every lift index $j$.
\end{enumerate}

\paragraph{Proof.}
Since $A\equiv1\pmod3$, $C\equiv3\pmod9$. Also $C$ is odd and $C\equiv3\pmod{q-1}$. Thus $3\nmid H$ and $\gcd(H,q-1)=1$. The moduli in \eqref{eq:label-selection} are pairwise coprime divisors of $R$. Multiplication of $k$ by a unit modulo $R$ preserves every condition. This proves independence of $g$, invariance under Frobenius, and invariance under any cyclotomic relabeling of the roots.

Multiplication of $z$ by $h\in\mu_N$ adds a multiple of $R/N$ to $k$. As $N\mid A-1$, the factor $C=3H$ divides $R/N$. For each $p\ne3$ dividing $N$, the integer $R/N$ is divisible by $(q-1)/N$, hence by $N$ and $p^6$. The selected class is therefore stable under every required scaling.

Write $k'$ for the exponent of $hz$. An element of $\mathbb F_{q^r}$ has degree smaller than $r$ exactly when it belongs to $\mathbb F_A$, since $r$ is a power of $3$. Applied to $(hz)^{\ell^u}$ this is exactly the divisibility $C\mid k'\ell^u$. When $\ell\ne3$ the factor $3$ prevents it. When $\ell=3$ the factor $H$ prevents it. The case $u=0$ gives degree $r$, hence an irreducible, squarefree monic polynomial with nonzero constant term. The general case gives the asserted orbit sizes.

The norm of $hz$ is $(g^e)^{k'}$, where $g^e$ generates $\mathbb F_q^{\times}$. All conjugates have this norm. At $p=3$ its order has the full $3$-part of $q-1$. At each $p\ne3$ dividing $N$, put $v=v_p(N)\ge6$. We have $v_p(q-1)\ge2v$ and $v_p(k')\le5$, so the order has $p$-adic valuation at least $2v-5\ge v$. This proves the divisibility for $y$.

The Chinese remainder theorem counts the permitted exponents by multiplying the proportions in \eqref{eq:label-count}. Every selected polynomial has $r$ roots. For the lower bound use $H\ge127$, $q^r\ge19^3$, and
\[
 \prod_{\substack{p\mid N\\p\ne3}}(1-p^{-6})
 \ge 1-\sum_{n=2}^\infty n^{-6}
 \ge 1-\frac1{64}-\int_2^\infty x^{-6}\,dx=\frac{313}{320}.
\]
Multiplying $(1-19^{-3})(2/3)(126/127)(313/320)>1/2$ gives the stated bound. Every primitive root meets \eqref{eq:label-selection}.

\paragraph{Passage to the projector.}
Keep the source's allowed-polynomial test in every coordinate. At one coordinate above each level prime require a selected polynomial. Its degree $r>1$ ensures it is allowed. Among all $q^r$ monic polynomials, at most a fraction $1-1/(2r)$ fails selection. With the source's uniform lower weight $w_0=(1-r/q)/q^r$, the total uniform weight of all $f$-tuples with no witness is at most $(1-1/(2r))^f\le e^{-f/(2r)}$. The existing allowed-label estimate still costs $f\epsilon$, where $\epsilon=(1+N+2r)/q$. Comparing the uniform weights to the actual regular-representation weights on the selected tuples gives
\[
 \operatorname{Re}[(1-e_{p_i})(1)]\le f\epsilon+e^{-f/(2r)}.
\]
Here $(1-e_{p_i})(1)$ denotes the identity coefficient in the group algebra. The second term is at most $b^{-9}$ for $f\ge18r\log b$. The two original numerical bounds on $f\epsilon+b^{-9}$ and $r<q$ still apply. On the Galois side the two conclusions of the finite-field theorem give every field of the source's \texttt{IsCommonRootPower}. In particular distinctness after raising to arbitrary $\ell$-powers yields unit differences after the tame-character embedding at $\ell$, exactly as in its existing matrix-span argument. One witness above each level prime ensures an available witness away from any fixed coefficient prime. Rank-one lift settings use their own label $X-1$ and their existing common-power lemma.

\paragraph{A strict enlargement.}
For $q=19$, $r=3$, $N=3$ the theorem's local hypotheses hold, with $R=6858$ and $H=127$. It selects $1512$ polynomials, compared with $\varphi(6858)/3=756$ primitive ones. In $\mathbb F_{19}[t]/(t^3-2)$ the element $g=1+t$ has order $6858$. The element $g^2$ has order $3429$ and minimal polynomial $X^3+16X^2+10X+10$. Its norm is $9$ of order $9$. This is a finite local example. The large-order parameter of the global proof is covered by the uniform argument above.

\paragraph{Lean integration.}
Change the witness in the \src{New/UpperBound/Objects/LabelsAndIndexTypes}{149}{good-tuple predicate} and propagate it to the family above both primes. Replace the \src{New/Spanning/MainTerm/BadTypes}{101}{totient count} and the \src{New/Spanning/MainTerm/BadTypes}{162}{nonprimitive mass bound} by \eqref{eq:label-count} and its exponential consequence. The \src{New/Spanning/SecondMoment/SpanningTheorem}{221}{identity-coefficient estimate} then uses the new threshold. In the \src{New/Spanning/Multiplicities/EqualMultiplicities}{249}{trace comparison}, the witness supplies irreducibility. Carry the new predicate into the high-rank branch of the \src{New/Congruences/GaloisRepresentations/SelmerHypotheses}{583}{lift setting} and its \src{New/Congruences/GaloisRepresentations/AttachedRepresentations}{81}{construction from good types}. The new norm and orbit arguments supply the \src{New/Congruences/GaloisRepresentations/RowsAndColumnsFromLabels}{330}{common-power structure} required by the \src{New/Congruences/GaloisRepresentations/RowsAndColumnsFromLabels}{548}{main-summand theorem}. The \src{New/Congruences/GaloisRepresentations/RowsAndColumnsFromLabels}{410}{rank-one construction} has its own witness. Enlarging the projector changes the spaces on both sides of the comparison. Their dimensions and multiplicities are related by applying the revised trace comparison to that same enlarged family.

\section{An explicit set of sparse infinite places}
\label{sec:paraboloid}

The derivative construction uses many places of large weight and requires that no three
of them occur at prescribed steps of a short Frobenius orbit. The current construction
in \path{New/Reduction/Data/Existence.lean} selects a largest sphere in a box of
integer digit vectors. One can instead list the exponents directly. The last two
digit positions record the squared length of the preceding digits. This gives an
exact cardinality and a membership test by integer division.

\paragraph{Finite construction.}
Take integers $n\ge3$, $H\ge1$, $L\ge2$, write $k=n-2$ and assume
$k\le L$. Put $B=LH$ and take any integer $o\ge B^n$. For
$u=(u_0,\ldots,u_{k-1})$ with $0\le u_i<H$, set
\[
  D(u)=\sum_{i=0}^{k-1}u_i B^i,\qquad
  Q(u)=\sum_{i=0}^{k-1}u_i^2,\qquad
  E(u)=D(u)+B^kQ(u).
\]
The set $S=\{E(u)\}$ lies in $[0,o)$ and has exactly $H^{n-2}$ elements.
For all positive integers $a,b$ with $a+b\le L$ it satisfies
\[
 (a+b)E(v)\equiv aE(u)+bE(w)\pmod o
        \quad\Longrightarrow\quad u=v=w .
\]
Indeed, each low digit on either side is at most $L(H-1)<B$.
The high block on either side is an integer between zero and
$Lk(H-1)^2<B^2$. Thus both sides are below $B^{k+2}\le o$ and are equal
as integers. Division by $B^k$ gives equality of the high blocks.
The remainders have unique base-$B$ digits, so
$(a+b)v=au+bw$ and $(a+b)Q(v)=aQ(u)+bQ(w)$. Substitution in the identity
\[
  (a+b)\bigl(aQ(u)+bQ(w)\bigr)-Q(au+bw)
        =ab\sum_{i=0}^{k-1}(u_i-w_i)^2
\]
makes its left side zero. Hence $u=w$ and then $v=u$.
Reduction of $E(u)$ modulo $B^k$ recovers $u$, which proves the cardinality.
Given an exponent, extracting those digits and comparing their squared length
with its quotient by $B^k$ decides membership. The cardinality is $n$ times
the lower bound $H^n/(nH^2)$ used in the existing proof. The comparison is with
that proved lower bound; a particular sphere can have more points.

\paragraph{Insertion at the existing parameters.}
Write $q_0$ for the auxiliary modulus, $s$ for the tower exponent,
$\ell_0=3^{\lfloor\log_3s\rfloor+4}$ for $s\ge1$, and
$q_1=q_0\cdot3\ell_0\,2^s$. Set $L=2\ell_0$ and $o=2^s$.
The digit side $H$ from the \exis{244}{sparse-set construction} already satisfies
$(2\ell_0 H)^n\le2^s$. For each fixed modulus, auxiliary prime and
$0<\varepsilon<1/100$, the proof of
the \exis{647}{large-scale existence lemma} chooses $n$ once;
its strict lower bound on $n$ permits $n\ge3$. The chosen ranks satisfy
$2^s=m^{22}\ge x^{11}$, so $s\to\infty$ and eventually
$n-2\le2\ell_0$. This verifies the one extra restriction for every
sufficiently large cutoff, with the existing order of quantifiers.

The places represented by $-g^{E(u)}\bmod q_1$, where $g=1+3\ell_0q_0$,
are in the fibre over $-1\bmod q_0$. The source exhibits a power-of-two
divisor of $q_1$ on which $g$ has exact order $o$. Thus distinct exponents give
distinct places. For an odd prime $p$, two forbidden ratios
$\pm p^i,\pm p^j$ with $0<i<j<\ell_0$ yield
\[
  2j E(v)\equiv 2(j-i)E(u)+2iE(w)\pmod{2^s}.
\]
The coefficient sum is $2j<L$. All three exponents coincide. The source's
order bound then forces $2^s\le(p+1)2i$, contrary to the admissible upper
bound $2\ell_0(P+1)<2^s$ for all tested primes $p\le P$.
For $p=2$, all selected places have the same reduction modulo $3\ell_0$,
and the order of $4$ modulo $3\ell_0$ is $\ell_0$. This rules out the first
ratio already. The exact cardinality of the explicit set therefore supplies
the same downstream bound for the number of places.

The formal change belongs at the \exis{181}{choice of progression-free exponents}. Carry $n\ge3$ and $n-2\le2\ell_0$ through the specialized \exis{244}{sparse-set} and \exis{488}{rank-dependent} lemmas to the \exis{898}{eventual construction}. The general sphere lemma still covers other parameter regimes. On the branch used for RH, the explicit formula supplies the set at every sufficiently large cutoff. The Frobenius order estimates, the prime windows, and the derivative divisibility estimates are still inputs.

\section{Comparison with the proposal archive and verification}\label{sec:verify}
The comparison uses ProveIt commit \href{https://github.com/VladimirReshetnikov/ProveIt/commit/89d12dd40424f1217fd818d84877b3da52396a3b}{\texttt{89d12dd4}}. Its \src{Main}{5}{readable Lean project} has Git tree \texttt{2e133df6}, identical to the readable project in upstream commit \href{https://github.com/lsernau/grh-lean/tree/72fead102f663d5c7a407964a13dc0cf9d977207/grh}{\texttt{72fead10}}. The \docref{guide}{human-readable guide} traces the digit construction in Section 8, the finite-field labels in Section 10, and the local operator assembly in Section 14. These are the three passages for which the present note supplies substitutions.

\begin{center}\small
\begin{tabularx}{\linewidth}{@{}p{.30\linewidth}X@{}}\toprule
Prior proposal, linked to source & Substitutions already proposed\\\midrule
\docref{proposals/centered_theta_lift_simplification}{Centered theta lifts} & Centered rank increment, ramified-source spanning, adaptive completion of a nearly full basis, and a random-product variant.\\[3pt]
\docref{proposals/proof_simplifications}{Three shorter steps} & Two-polynomial and Fourier estimates for the auxiliary prime, a Perron argument at the far cutoff, and a stronger spanning estimate.\\[3pt]
\docref{proposals/rh-simplifications}{Simplifications of RH} & Rank recurrence on the ramified source, simultaneous torsion interpolation for choosing twists, and a quadratic-large-sieve selector.\\[3pt]
\docref{proposals/rh_sparse_theta_proof}{Sparse cutoffs} & Reciprocal-convex interpolation, sparse or averaged cutoffs, and a revised quadratic-twist selection argument.\\[3pt]
\docref{proposals/shorter-theta-lift-proof}{Shortening the theta-lift proof} & Smooth blocks and a large-sieve route for choosing the quadratic field, ramified-source spanning, and endpoint interpolation.\\\bottomrule
\end{tabularx}\end{center}

In particular, torsion interpolation chooses a character before constructing the one-generator operators. Section~\ref{sec:flag} combines operators after that arithmetic work is finished. Adaptive basis completion estimates the cost of missing target vectors in the covolume comparison. The flag used here is inside one local lift space and its successive quotients all have dimension one. In Section~\ref{sec:denser-labels} the count concerns local polynomial types before the spanning argument starts. In Section~\ref{sec:paraboloid} the sparse objects are infinite places of a CM field at one fixed cutoff. The sparse-cutoff proposal varies the cutoff itself. These distinctions allow the three changes to be evaluated even after the existing proposals have been implemented.

\paragraph{Verification of the baseline.}
The \href{https://math.mit.edu/~drew/shortgrh/}{independent verification record of Andrew Sutherland} reports a full build, kernel replay of every module, and a successful paranoid Comparator run for the original readable edition. Its toolchain is Lean 4.35.0 with Mathlib revision \texttt{56a00585}. The equality of Git trees identifies which source those checks cover. This report does not reproduce that build.

\paragraph{Verification of the replacements.}
The attached standalone Lean files use Lean 4.32.0 and its standard library. \path{formal/IntegralPivots.lean} proves the existence of each successive pivot, the extension of the flag by a normalized integral vector, and the equivalence between coordinate integrality and integrality of its expansion. \path{formal/FlagRows.lean} proves in all finite dimensions that selecting the prescribed rows preserves clearing and gives the stated determinant. It defines the determinant by cofactor expansion and proves the triangular formula for that definition. \path{formal/DigitRigidity.lean} proves the modular rigidity of the explicit code for all admissible finite parameters, including the range bound, decoding, and injectivity. \path{formal/LabelArithmetic.lean} proves the integer obstruction to orbit collapse, invariance of the residue tests under the corresponding shifts, and the prime-power divisibility used in the norm estimate. The new modules are separate from the pinned project. The paper supplies the identification of the integral evaluation lattice with adelic forms, the finite-field orbit count and its application to the projector, and the transfer of the digit code to Frobenius orbits of places. Integration and measurement of the revised project remain open.

The three exact-arithmetic programs exercise the pivot construction and clearing on rational models at several prime valuations, the orbit and norm tests in finite fields, and the digit and signed-place conditions. The archive includes their outputs, the Lean axiom reports, and a source index with immutable links. Its checker verifies the indicated lines and Git blob hashes against a checkout. The finite examples supplement the uniform proofs.

\paragraph{Which change to implement first.}
The row construction can be inserted between the existing one-generator and determinant lemmas, then reused for the compressed determinant. The digit code can be integrated separately by strengthening only the eventual hypotheses on its fixed dimension. The projector change has the largest quantitative effect among the three local estimates and the widest integration surface. The test for success there is to build both the trace-comparison and Galois branches for the enlarged projector, as well as the final bad-type estimate. No global improvement to the required field degree or to the running time follows solely from the new local threshold.

\paragraph{Further questions.}
The individual bounds $B_j$ can be summed before taking a worst case. Do the Selmer estimates yield a smaller sum for the families actually used in the proof? The finite-field criterion also leaves room for an optimum. Roots whose exponent is divisible by $3$ can still retain their full orbit after every prime-power map if the remaining order in the quotient by $\mathbb F_{q^{r/3}}^\times$ has two distinct prime divisors. Counting these additional labels would require information about $H$, which the uniform estimate here avoids.

\end{document}
